\section{Methodology}

This study follows a machine learning pipeline composed of exploratory data analysis,
preprocessing, model training and validation, and interpretation of results.
Each step is described in detail below.

\subsection{Exploratory Data Analysis (EDA)}

\todo{
\begin{itemize}
    \item Describe the main characteristics of the dataset, including the number of instances, number
        and types of attributes, and the definition of the target variable.
    \item Present the distribution of the target attribute and discuss possible class imbalance issues.
    \item Report the presence of patterns, anomalies, or inconsistencies observed in the dataset, such
        as outliers, missing values, duplicated records, or irregular values.
    \item Present visualizations and statistical summaries used to support the exploratory analysis
        (e.g., histograms, boxplots, correlation matrices, descriptive statistics).
    \item Discuss relationships between attributes and identify possible redundant or highly correlated features.
    \item Highlight attributes or instances that were identified as potentially irrelevant, redundant, or problematic.
    \item Summarize the main findings obtained from the exploratory analysis.
    \item Based on these findings, describe the data quality issues identified and motivate the preprocessing steps
        required for data preparation.
\end{itemize}
}

\subsection{Data Preprocessing}

\todo{
\begin{itemize}
    \item Describe the preprocessing steps applied to the dataset based on the issues identified during the
        exploratory data analysis.
    \item Explain the criteria used to remove or modify attributes or instances considered
        irrelevant, redundant, or inconsistent.
    \item Describe how missing values were handled and justify the chosen strategy.
    \item Report how outliers were treated and explain the reasoning behind the adopted approach.
    \item Describe any normalization or standardization procedures applied to numerical attributes.
    \item Explain how categorical variables were encoded and justify the selected encoding method.
    \item Describe any techniques used to address class imbalance, when applicable.
    \item Report the use of dimensionality reduction or feature selection methods, if applied,
        including the criteria used.
    \item Explain how preprocessing steps were incorporated into the modeling workflow.
    \item If no preprocessing was required, explicitly state this and justify the decision based on the data analysis.
\end{itemize}
}

\subsection{Definition of the Approach, Algorithms, and Evaluation Strategy}

\todo{
\begin{itemize}
    \item Describe the overall machine learning approach adopted to address the problem, including
        the type of learning task (e.g., classification or regression).
    \item Present the set of supervised learning algorithms selected for the study and justify their
        inclusion based on their characteristics and suitability for the problem.
    \item Highlight the diversity of the selected algorithms in terms of modeling strategies or inductive biases.
    \item Describe the performance evaluation metrics used to assess model performance, explaining
        their relevance to the problem context.
    \item Indicate the primary metric used as the main criterion for model selection and optimization,
        and justify this choice.
    \item Describe the strategy used to divide the dataset into training, validation, and testing subsets.
    \item Explain the evaluation methodology adopted (e.g., holdout validation or cross-validation),
        including how it contributes to reliable model assessment.
    \item Justify the methodological choices made in this stage based on best practices or characteristics
        of the dataset and problem.
\end{itemize}
}

\subsection{Model Training, Hyperparameter Tuning, and Validation}

\todo{
\begin{itemize}
    \item Describe how the selected models were trained according to the previously defined experimental setup.
    \item Report whether an initial comparison of algorithms (e.g., spot-checking) was performed and describe its
        purpose in guiding the selection of promising models.
    \item Describe the hyperparameter tuning procedures applied to the selected models, including the search strategies used.
    \item Explain how data separation was maintained throughout the training and validation process to avoid data leakage.
    \item Describe how consistency across experiments was ensured, including the use of identical data partitions when
        comparing different algorithms.
    \item Report the use of repeated experiments or multiple runs, when applicable, to improve the reliability of
        the obtained results.
    \item Describe the steps taken to ensure reproducibility of the experiments, such as the use of fixed random seeds.
    \item Summarize how model performance was recorded and organized for subsequent analysis.
\end{itemize}
}

\subsection{Model Interpretation and Analysis}

After model selection, post-hoc explainability methods were applied to the
selected CKD classifier. This interpretation stage was not used to train or
select the model; instead, it was used to analyze the behavior of the selected
model after the predictive modeling workflow had been completed. The
explainability analysis combined model-specific feature importance, permutation
feature importance, SHAP-based summaries, SHAP waterfall plots, and LIME local
explanations \cite{fisher2019modelreliance,lundberg2017shap,ribeiro2016lime}.
Because the selected classifier was a Random Forest model
\cite{breiman2001randomforests}, native feature importance was also reported as
a model-specific global explanation.

Global explanations were used to identify features most associated with the
model predictions across the test set, while local explanations were used to
inspect representative individual predictions. SHAP values were interpreted
with respect to the configured explained class, \texttt{ckd}; positive SHAP
values indicate increased support for that class, whereas negative SHAP values
indicate reduced support for that class. These explanations were treated as
model-dependent interpretability evidence and not as causal clinical
conclusions \cite{molnar2022interpretable,doshivelez2017interpretable}.
